\documentclass[exercice]{thedocument}%
\startdatakeys
\source{}
\BO{2026}
\cycle{4}
\competencesDuSocle{CA}
\objectifsApprentissage{C4A2bg,C4D2ad,C4A3bc}
\automatismes{C4D2aa,C4A2bd}
\enddatakeys
\begin{document}%
Un confiseur veut confectionner des sachets contenant des sucettes
à la menthe et des sucettes au citron. Les sucettes à la menthe lui reviennent
à \<n1;{\euro} l'une et les sucettes au citron lui reviennent à \<n2;{\euro}
l'une. Dans chaque sachet, le confiseur veut qu'il y ait \<n3; sucettes et
il souhaite que le prix de revient d'un sachet ne dépasse pas \<n4;\euro.
\begin{itemize}
    \item Déterminer toutes les compositions de sachets possibles.
\end{itemize}
\end{document}
